% ======================================================================
%  PART 8 -- Forecasting
% ======================================================================
\section{Forecasting}
\label{sec:forecasting}

The purpose of time series analysis is often to \emph{forecast} future values
from the data observed up to now. The recipe has three steps: (1) assume the
form of the model; (2) estimate its parameters; (3) use the assumed structure
and the estimated parameters to form predictions. In this part we take the
model structure and the parameters as \textbf{known} and predict $h$ lags
ahead, i.e.\ we want $X_{n+h}$ from $X_1,\dots,X_n$. The integer $h>0$ is the
\textbf{lead} (or \textbf{horizon}) and $n$ is the \textbf{origin}. The central
theme is that, for the second-order theory we use, the optimal predictor depends
\emph{only} on the covariance structure $\{\acvs_\tau\}$, and for causal ARMA
models it is computed cleanly by setting future innovations to zero.

% ----------------------------------------------------------------------
\subsection{Definitions}

\begin{definition}{Prediction mean square error}{p8-pmse}
Let $g(X_1,\dots,X_n)$ be any estimator (predictor) of $X_{n+h}$ based on the
observations $X_1,\dots,X_n$. Its performance is measured by the
\textbf{prediction mean square error}
\[
P^{\,n}_{n+h}\;:=\;\E\!\left[\big(X_{n+h}-g(X_1,\dots,X_n)\big)^2\right].
\]
Smaller $P^{\,n}_{n+h}$ means a better forecast; minimising it over a class of
predictors defines the ``best'' predictor in that class.
\end{definition}

\begin{definition}{Best linear predictor}{p8-blp}
The \textbf{best linear predictor} of $X_{n+h}$ based on $X_1,\dots,X_n$ is the
linear function
\[
X^{\,n}_{n+h}\;=\;P_{X_1,\dots,X_n}(X_{n+h})\;=\;\sum_{j=1}^{n}\beta_j X_j
\]
that minimises the prediction mean square error among all linear combinations of
$X_1,\dots,X_n$. We write $P_{X_1,\dots,X_n}(\cdot)$ for the projection onto the
linear span of the data.
\begin{itemize}
  \item For \textbf{Gaussian} time series the conditional expectation
        $\E[X_{n+h}\mid X_1,\dots,X_n]$ is itself linear in the data, so the best
        linear predictor coincides with the (overall) best predictor.
  \item In general the conditional expectation is a \emph{nonlinear} function of
        the data; for tractability one restricts attention to linear predictors,
        which need only second-order information.
\end{itemize}
\end{definition}

\begin{definition}{$h$-step forecast error and its variance}{p8-ferr}
For the best linear predictor $X^{\,n}_{n+h}$ define the \textbf{$h$-step ahead
forecast error}
\[
e_n(h)\;:=\;X_{n+h}-X^{\,n}_{n+h}.
\]
It is unbiased, $\E[e_n(h)]=0$, and its variance equals the prediction MSE,
$\Var\big(e_n(h)\big)=P^{\,n}_{n+h}$. For a causal ARMA process with
$MA(\infty)$ weights $\{\psi_j\}$ one has
$e_n(h)=\sum_{j=0}^{h-1}\psi_j\eps_{n+h-j}$ and
$\Var\big(e_n(h)\big)=\sigma^2\sum_{j=0}^{h-1}\psi_j^2$.
\end{definition}

\begin{definition}{Sources of prediction uncertainty}{p8-srcunc}
There are \textbf{three} components of uncertainty in a forecast:
\begin{enumerate}
  \item \textbf{model uncertainty} --- the assumed model form may be wrong;
  \item \textbf{estimation uncertainty} --- the parameters are estimated, not known;
  \item \textbf{innovation uncertainty} --- the inherent randomness of the future
        innovations $\eps_{n+1},\dots,\eps_{n+h}$.
\end{enumerate}
Throughout we assume the model is \emph{known} and the parameters are estimated
\emph{without error}, which removes the first two sources and leaves only
innovation uncertainty. (Bootstrap methods can incorporate the first two, but
are not covered here.)
\end{definition}

\begin{definition}{Gaussian prediction interval}{p8-pint}
For a \textbf{Gaussian} process the $h$-step error $e_n(h)$ is normal with mean
$0$ and variance $\Var(e_n(h))$, so a $(1-\alpha)$ \textbf{prediction interval}
for $X_{n+h}$ based on $X_1,\dots,X_n$ is
\[
X^{\,n}_{n+h}\;\pm\;z_{1-\alpha/2}\,\sqrt{\Var\big(e_n(h)\big)},
\]
where $z_{1-\alpha/2}$ is the $(1-\alpha/2)$ quantile of the standard normal
distribution (e.g.\ $z_{0.975}\approx1.96$ for a $95\%$ interval).
\end{definition}

\begin{intuition}
A point forecast alone is almost useless without a measure of its spread. The
interval widens with $h$ because more unknown innovations
$\eps_{n+1},\dots,\eps_{n+h}$ accumulate in the error. For a stationary ARMA
model the width saturates: as $h\to\infty$, $\Var(e_n(h))\to\acvs_0$, i.e.\ the
interval converges to ``mean $\pm$ a multiple of the marginal standard
deviation'' --- forecasting far ahead is no better than quoting the
unconditional distribution. For integrated (ARIMA) processes the interval grows
\emph{without bound}.
\end{intuition}

% ----------------------------------------------------------------------
\subsection{Theorems \& Propositions}

\begin{lemma}{Conditional expectation minimises the prediction MSE}{p8-condexp}
The conditional expectation
\[
g(X_1,\dots,X_n)=\E\big[X_{n+h}\mid X_1,\dots,X_n\big]
\]
minimises the prediction mean square error for $X_{n+h}$ based on
$X_1,\dots,X_n$.\\[3pt]
\textit{Proof idea:} For any real random variable $Y$ with $\mu=\E[Y]$ and any
constant $c$,
\[
\MSE(c)=\E\big[(Y-c)^2\big]=\E\big[(Y-\mu+\mu-c)^2\big]
=\E\big[(Y-\mu)^2\big]+(\mu-c)^2=\Var(Y)+(\mu-c)^2,
\]
which is minimised at $c=\mu=\E[Y]$ (the cross term vanishes because
$\E[Y-\mu]=0$). Applying this \emph{conditionally on} $X_1,\dots,X_n$ and using
the tower property
$\E[(X_{n+h}-g)^2]=\E\big[\E[(X_{n+h}-g)^2\mid X_1,\dots,X_n]\big]$, the inner
MSE is minimised pointwise by taking $g=\E[X_{n+h}\mid X_1,\dots,X_n]$.\\[2pt]
\textit{Why:} this is the theoretical gold standard for forecasting. For
Gaussian processes the conditional expectation is linear, so the best predictor
\emph{is} the best linear predictor; in general it motivates restricting to
linear predictors, which only need the ACVS.
\end{lemma}

\begin{example}{Gaussian AR(1) predictor --- mean reversion}{p8-ar1}
Let $X_t=\phi X_{t-1}+\eps_t$ with $\eps_t\iid\mathcal N(0,\sigma^2)$ and
$\abs\phi<1$ (stationary). Because the process is Gaussian, the best linear
predictor equals the conditional expectation, and iterating the recursion,
\[
X^{\,n}_{n+h}=\E[X_{n+h}\mid X_1,\dots,X_n]
=\E[\phi X_{n+h-1}+\eps_{n+h}\mid X_1,\dots,X_n]
=\phi^{\,h}X_n .
\]
Since $\abs\phi<1$, $X^{\,n}_{n+h}=\phi^{\,h}X_n\to0$ (the process mean) as
$h\to\infty$: the forecast \textbf{reverts to the mean}. This is a general
property of stationary ARMA models. (A worked figure in the notes uses
$\phi=0.8$, $\sigma^2=1$.)
\end{example}

\begin{example}{Gaussian MA(1) predictor}{p8-ma1}
Let $X_t=\mu+\eps_t-\theta\eps_{t-1}$ with $\eps_t\iid\mathcal N(0,\sigma^2)$.
For $h>0$,
\[
X^{\,n}_{n+h}=\mu+\E[\eps_{n+h}\mid X_1,\dots,X_n]-\theta\,\E[\eps_{n+h-1}\mid X_1,\dots,X_n]
=\mu-\theta\,\E[\eps_{n+h-1}\mid X_1,\dots,X_n],
\]
because $\eps_t$ for $t>n$ is uncorrelated with $X_1,\dots,X_n$ (so its
conditional expectation is $0$). Hence:
\begin{itemize}
  \item $h=1$: $\;X^{\,n}_{n+1}=\mu-\theta\,\hat\eps_n$ (the one observable
        innovation $\eps_n$ remains);
  \item $h>1$: $\;X^{\,n}_{n+h}=\mu$ --- the forecast hits the mean immediately,
        consistent with an MA(1) having memory only of length $1$.
\end{itemize}
\end{example}

\begin{theorem}{Prediction equations (zero-mean stationary process)}{p8-predeq}
For a zero-mean stationary process $\{X_t\}$, the best linear predictor
$X^{\,n}_{n+h}=\sum_{j=1}^n\beta_j X_j$ is found by solving for
$\beta_1,\dots,\beta_n$ the \textbf{prediction equations}
\[
\E\!\big[(X_{n+h}-X^{\,n}_{n+h})\,X_k\big]=0,\qquad k=1,\dots,n,
\]
i.e.\ the forecast error is orthogonal to every observation. Equivalently, in
matrix form $\Gamma_n\beta=\acvs_{[h]}$:
\[
\begin{pmatrix}
\acvs_0 & \acvs_1 & \cdots & \acvs_{n-1}\\
\acvs_1 & \acvs_0 & \cdots & \acvs_{n-2}\\
\vdots & \vdots & \ddots & \vdots\\
\acvs_{n-1} & \acvs_{n-2} & \cdots & \acvs_0
\end{pmatrix}
\begin{pmatrix}\beta_1\\ \beta_2\\ \vdots\\ \beta_n\end{pmatrix}
=\begin{pmatrix}\acvs_{n+h-1}\\ \acvs_{n+h-2}\\ \vdots\\ \acvs_{h}\end{pmatrix}.
\]
\textit{Proof idea:} the best linear predictor is the orthogonal projection of
$X_{n+h}$ onto $\mathrm{span}\{X_1,\dots,X_n\}$; the residual
$X_{n+h}-X^{\,n}_{n+h}$ must be orthogonal (uncorrelated) to each $X_k$.
Expanding $\E[(X_{n+h}-\sum_j\beta_jX_j)X_k]=0$ gives
$\acvs_{n+h-k}=\sum_{j=1}^n\beta_j\acvs_{j-k}$, which is row $k$ of
$\Gamma_n\beta=\acvs_{[h]}$.\\[2pt]
\textit{Why:} this is the master equation for linear prediction of \emph{any}
stationary process --- it shows the optimal forecast uses only the ACVS.
\end{theorem}

\begin{corollary}{Matrix solution and prediction MSE}{p8-matsol}
If $\Gamma_n$ is invertible, then with $X=(X_1,\dots,X_n)\Tr$ the best linear
predictor and its MSE are
\[
X^{\,n}_{n+h}=\acvs_{[h]}\Tr\,\Gamma_n^{-1}\,X,
\qquad
P^{\,n}_{n+h}=\acvs_0-\acvs_{[h]}\Tr\,\Gamma_n^{-1}\,\acvs_{[h]} .
\]
\textit{Proof idea:} solve $\Gamma_n\beta=\acvs_{[h]}$ for
$\beta=\Gamma_n^{-1}\acvs_{[h]}$ and substitute into
$X^{\,n}_{n+h}=\beta\Tr X$ and into
$P^{\,n}_{n+h}=\acvs_0-2\beta\Tr\acvs_{[h]}+\beta\Tr\Gamma_n\beta$; the last
two terms collapse using $\Gamma_n\beta=\acvs_{[h]}$.\\[2pt]
\textit{Why:} gives the closed-form forecast for short series.
\begin{itemize}
  \item $\Gamma_n$ is invertible for \emph{invertible} ARMA models (and not in
        all other cases), but the best prediction $X^{\,n}_{n+h}$ is
        \textbf{always unique} even when $\Gamma_n$ is singular.
  \item Computing $\acvs_{[h]}\Tr\Gamma_n^{-1}X$ is inefficient for large $n$
        because the $n\times n$ matrix must be inverted; the Durbin--Levinson
        and innovations algorithms avoid this (Shumway \& Stoffer, \S3.5).
\end{itemize}
\end{corollary}

\begin{theorem}{Prediction equations with non-zero mean (and intercept $\beta_0$)}{p8-predeqmean}
For a stationary process with $\mu=\E[X_t]$, write
$X^{\,n}_{n+h}=\beta_0+\sum_{j=1}^n\beta_j X_j$. The coefficients solve the
prediction equations $\E[(X_{n+h}-X^{\,n}_{n+h})X_k]=0$ for $k=0,1,\dots,n$
(with the convention $X_0\equiv1$), which give
\[
\Gamma_n\beta=\acvs_{[h]},
\qquad
\beta_0=\mu\big(1-\beta\Tr\mathbf 1_n\big),
\]
and, when $\Gamma_n$ is invertible,
\[
X^{\,n}_{n+h}=\mu+\acvs_{[h]}\Tr\,\Gamma_n^{-1}(X-\mu\mathbf 1_n),
\qquad
P^{\,n}_{n+h}=\acvs_0-\acvs_{[h]}\Tr\,\Gamma_n^{-1}\,\acvs_{[h]} .
\]
\textit{Proof idea:} the predictor has the form
$X^{\,n}_{n+h}=\mu+\beta\Tr(X-\mu\mathbf 1_n)$, so centring removes the mean and
reduces the problem to the zero-mean case --- there is no loss of generality in
assuming $\mu=0$, where $\beta_0=0$. The MSE is \emph{identical} to the
zero-mean case because the constant shift does not affect the error variance.\\[2pt]
\textit{Why:} real data are rarely centred; this is the version actually used in
practice. (Proved as Exercise~\ref{exo:p8-ex2}.)
\end{theorem}

\begin{theorem}{Best linear predictor of a causal ARMA via $\psi$-weights}{p8-armapred}
For a causal ARMA process with $MA(\infty)$ representation
$X_t=\sum_{j=0}^\infty\psi_j\eps_{t-j}$, the best linear predictor of $X_{n+h}$
based on $X_1,\dots,X_n$ is
\[
X^{\,n}_{n+h}=\sum_{j=h}^{\infty}\psi_j\,\eps_{n+h-j}
=\psi_h\eps_n+\psi_{h+1}\eps_{n-1}+\cdots,
\]
with prediction mean square error
\[
P^{\,n}_{n+h}=\sigma^2\sum_{j=0}^{h-1}\psi_j^2 .
\]
\textit{Proof idea:} write the predictor as a linear filter of the past
innovations, $X^{\,n}_{n+h}=\sum_{j=0}^\infty c_j\eps_{n-j}$. Then
\[
\E\big[(X_{n+h}-X^{\,n}_{n+h})^2\big]
=\sigma^2\Big(\sum_{j=0}^{h-1}\psi_j^2+\sum_{j=h}^\infty(\psi_j-c_{j-h})^2\Big),
\]
which is minimised by matching $c_{j}=\psi_{j+h}$, killing the second sum and
leaving $\sigma^2\sum_{j=0}^{h-1}\psi_j^2$. Practically:
\textbf{set the unrealised future innovations $\eps_{n+1},\dots,\eps_{n+h}$ to
zero} in $X_{n+h}=\sum_{j=0}^\infty\psi_j\eps_{n+h-j}$.\\[2pt]
\textit{Why:} this is the engine of all ARMA forecasting. It also shows that for
a causal ARMA, $\psi_j\to0$ exponentially so $X^{\,n}_{n+h}\to\mu$, and
$P^{\,n}_{n+h}=\sigma^2\sum_{j=0}^{h-1}\psi_j^2\uparrow\sigma^2\sum_{j=0}^\infty\psi_j^2=\acvs_0$
as $h\to\infty$. (Proved as Exercise~\ref{exo:p8-ex3}.)
\end{theorem}

\begin{intuition}
The two ARMA forecasting theorems (prediction equations vs.\ $\psi$-weights) are
the same projection viewed differently. The prediction equations project onto
$\mathrm{span}\{X_1,\dots,X_n\}$ using the ACVS; the $\psi$-weight formula
projects onto $\mathrm{span}\{\eps_s:s\le n\}$, the \emph{innovation} space.
Because for an invertible ARMA these two spans coincide, the answers agree --- but
the innovation view is far easier to compute: ``write the future, erase the
future innovations.''
\end{intuition}

\begin{proposition}{Limiting behaviour of ARMA forecasts}{p8-armalimit}
For a causal ARMA model with mean $\mu$,
$X^{\,n}_{n+h}=\mu+\sum_{j=h}^\infty\psi_j\eps_{n+h-j}$ with $\psi_j\to0$
exponentially fast, hence
\[
X^{\,n}_{n+h}\longrightarrow\mu,
\qquad
\Var\big(e_n(h)\big)=\sigma^2\sum_{j=0}^{h-1}\psi_j^2\longrightarrow\acvs_0,
\qquad h\to\infty .
\]
\textit{Proof idea:} immediate from the $\psi$-weight theorem and exponential
decay of $\psi_j$ (a consequence of causality, roots of $\Phi$ outside the unit
circle).\\[2pt]
\textit{Why:} quantifies ``the distant future is the unconditional
distribution'' --- the point forecast converges to the mean and the error
variance converges to the marginal variance $\acvs_0$.
\end{proposition}

\begin{proposition}{Forecasts for integrated (ARIMA) processes}{p8-arimapred}
Let $\{X_t\}$ be ARIMA$(p,d,q)$, $\Phi(B)(1-B)^dX_t=\Theta(B)\eps_t$. Then
$Z_t=\diff^{\,d}X_t$ is the ARMA$(p,q)$ process
$Z_t=\sum_{j=1}^p\phi_jZ_{t-j}+\eps_t-\sum_{j=1}^q\theta_j\eps_{t-j}$, and one
forecasts $X$ by undoing the differencing,
\[
\diff^{\,d}X^{\,n}_{n+h}=Z^{\,n}_{n+h}.
\]
For $d=1$ (so $X_{n+h}-X_{n+h-1}=Z_{n+h}$),
\[
X^{\,n}_{n+h}=X_n+\sum_{j=0}^{h-1}Z^{\,n}_{n+h-j},\qquad h\ge1 .
\]
If $\{Z_t\}$ has mean $\alpha\ne0$ then $Z^{\,n}_{n+h}\to\alpha$ and
$X^{\,n}_{n+h}\approx X_n+\alpha h$ for large $h$ (a deterministic drift line).\\[2pt]
\textit{Proof idea:} apply the linear projection operator $P_{X_1,\dots,X_n}$,
which is linear, to $\diff^{\,d}X_{n+h}=Z_{n+h}$; since differencing is a fixed
linear filter it commutes with the projection, giving
$\diff^{\,d}X^{\,n}_{n+h}=Z^{\,n}_{n+h}$. Telescoping the $d=1$ identity
$X_{n+h}=X_{n+h-1}+Z_{n+h}$ down to the last observation $X_n$ yields
$X^{\,n}_{n+h}=X_n+\sum_{j=0}^{h-1}Z^{\,n}_{n+h-j}$.\\[2pt]
\textit{Why:} reduces ARIMA forecasting to ARMA forecasting on the differenced
series, then ``re-integrates''.
\begin{itemize}
  \item For ARIMA models the prediction MSE \textbf{diverges}:
        $P^{\,n}_{n+h}\to\infty$ as $h\to\infty$ (no mean reversion --- the
        process is non-stationary).
  \item For SARIMA models with $d=D=1$ and $\alpha=0$ the forecast is linear and
        seasonal in $h$.
\end{itemize}
\end{proposition}

\begin{pitfall}
After expanding an ARIMA$(p,1,q)$ into the form
$X_t=(1+\phi_1)X_{t-1}-(\phi_1-\phi_2)X_{t-2}-\cdots$, the equation \emph{looks}
like an ARMA$(p{+}1,0,q)$, but it is \textbf{not} a stationary ARMA: its AR
coefficients do not satisfy the causality/stationarity conditions (there is a
unit root). Treat it only as a convenient recursion for forecasting, never as a
stationary model. This is exactly why $P^{\,n}_{n+h}\to\infty$.
\end{pitfall}

% ----------------------------------------------------------------------
\subsection{Key Techniques}

\begin{keytech}{Solving the prediction equations}{p8-ktpredeq}
To get the best linear predictor of a zero-mean stationary process:
\begin{enumerate}
  \item Build $\Gamma_n=[\acvs_{i-j}]_{i,j=1}^n$ (Toeplitz, from the ACVS).
  \item Build the right-hand side
        $\acvs_{[h]}=(\acvs_{n+h-1},\acvs_{n+h-2},\dots,\acvs_h)\Tr$.
  \item Solve $\Gamma_n\beta=\acvs_{[h]}$ for $\beta$, set
        $X^{\,n}_{n+h}=\sum_j\beta_jX_j=\acvs_{[h]}\Tr\Gamma_n^{-1}X$.
  \item Report the MSE
        $P^{\,n}_{n+h}=\acvs_0-\acvs_{[h]}\Tr\Gamma_n^{-1}\acvs_{[h]}$.
\end{enumerate}
For \textbf{non-zero mean} $\mu$, centre first: use $X-\mu\mathbf1_n$ in place of
$X$ and add $\mu$ back,
$X^{\,n}_{n+h}=\mu+\acvs_{[h]}\Tr\Gamma_n^{-1}(X-\mu\mathbf1_n)$; the MSE is
unchanged. \emph{Shortcut for AR(p):} for $h=1$ and $n\ge p$ the answer is just
$X^{\,n}_{n+1}=\sum_{i=1}^p\phi_iX_{n+1-i}$ --- no inversion needed (the next
technique explains why).
\end{keytech}

\begin{keytech}{Causal-ARMA forecast by zeroing future innovations}{p8-ktpsi}
The fast route for any causal ARMA:
\begin{enumerate}
  \item Write $X_{n+h}=\sum_{j=0}^\infty\psi_j\eps_{n+h-j}$ (the $MA(\infty)$
        form), or just the ARMA equation rearranged for $X_{n+h}$.
  \item Take conditional expectation given the data: every future innovation
        $\eps_{n+1},\dots,\eps_{n+h}$ becomes $0$; every past/present
        innovation $\eps_s$ ($s\le n$) stays (replace by its residual
        $\hat\eps_s$); every past observation $X_s$ ($s\le n$) stays; every
        future observation $X_s$ ($s>n$) is replaced by its forecast
        $X^{\,n}_s$.
  \item The MSE is $P^{\,n}_{n+h}=\sigma^2\sum_{j=0}^{h-1}\psi_j^2$ --- obtain the
        first $h$ $\psi$-weights from $\Theta(B)/\Phi(B)=\sum_j\psi_jB^j$ by
        matching powers of $B$.
\end{enumerate}
\end{keytech}

\begin{keytech}{Recursive truncated forecasting recipe (ARMA \& ARIMA)}{p8-ktrec}
The general-purpose hand-computation method. \textbf{Three steps:}
\begin{enumerate}
  \item \textbf{Isolate $y_t$:} expand the (ARI)MA equation so that $y_t$ is
        alone on the left and \emph{everything else} (past $y$'s, present/past
        $\eps$'s) is on the right.
  \item \textbf{Shift to $n+h$:} replace $t$ by $n+h$ throughout.
  \item \textbf{Substitute:} on the right-hand side replace
        \begin{itemize}
          \item future observations $y_{s}$ ($s>n$) by their forecasts
                $\hat y_{s\mid n}=X^{\,n}_s$;
          \item future errors $\eps_s$ ($s>n$) by $0$;
          \item past errors $\eps_s$ ($s\le n$) by the corresponding residuals
                $\hat\eps_s$ (and any $\eps_s$ with $s\le0$ by $0$).
        \end{itemize}
\end{enumerate}
Run $h=1,2,3,\dots$ in order; each forecast feeds the next. The residuals are
themselves built recursively (for $t=1,\dots,n$, and $0$ otherwise) via
\[
\hat\eps_t=\Big(\tilde X^{\,n}_t-\sum_{i=1}^p\phi_i\tilde X^{\,n}_{t-i}\Big)
+\sum_{k=1}^q\theta_k\hat\eps_{t-k}.
\]
The truncation (setting pre-sample terms to $0$) makes $\tilde X^{\,n}_t$ an
\emph{approximation}; the forecast error is estimated by
$\sigma^2\sum_{j=0}^{h-1}\psi_j^2$.
\end{keytech}

\begin{keytech}{Re-integrating ARIMA forecasts}{p8-ktint}
For ARIMA$(p,d,q)$: difference $d$ times to get the ARMA series
$Z_t=\diff^{\,d}X_t$, forecast $Z^{\,n}_{n+h}$ by the ARMA technique, then undo
the differencing. For $d=1$,
\[
X^{\,n}_{n+1}=X_n+Z^{\,n}_{n+1},\qquad
X^{\,n}_{n+h}=X_n+\sum_{j=0}^{h-1}Z^{\,n}_{n+h-j}\ (h\ge1).
\]
Equivalently, expand $\Phi(B)(1-B)^d$ as a single (unit-root) AR polynomial in
$X$ and apply the recursive recipe directly (often easier in practice). Remember
the MSE \emph{grows without bound}.
\end{keytech}

\begin{keytech}{Building Gaussian prediction intervals}{p8-ktpi}
\begin{enumerate}
  \item Get the point forecast $X^{\,n}_{n+h}$.
  \item Get the $\psi$-weights $\psi_0=1,\psi_1,\dots,\psi_{h-1}$ from
        $\Theta(B)/\Phi(B)$.
  \item Forecast-error variance: $\Var(e_n(h))=\sigma^2\sum_{j=0}^{h-1}\psi_j^2$.
  \item Interval: $X^{\,n}_{n+h}\pm z_{1-\alpha/2}\sqrt{\Var(e_n(h))}$ (use
        $1.96$ for $95\%$).
\end{enumerate}
Sanity check: for stationary ARMA the half-width increases in $h$ toward
$z_{1-\alpha/2}\sqrt{\acvs_0}$; for ARIMA it keeps growing.
\end{keytech}

% ----------------------------------------------------------------------
\subsection{Exercises}

\begin{exercise}{One-step AR(2) predictors from 1 and 2 observations \quad\textnormal{[Sheet 10]}}{p8-ex1}
Let $\eps_t$ be white noise. For the causal AR(2) process
$X_t=\phi_1X_{t-1}+\phi_2X_{t-2}+\eps_t$, show that the best one-step-ahead
linear predictors based on one observation $X_1$ and on two observations
$X_1,X_2$ are
\[
X^{\,1}_2=\frac{\acvs_1}{\acvs_0}X_1=\acf_1X_1,
\qquad
X^{\,2}_3=\phi_1X_2+\phi_2X_1 .
\]
\end{exercise}
\begin{solution}
\textbf{Based on $X_1$ (so $n=1$, $h=1$).} The prediction equation
$\Gamma_1\beta=\acvs_{[1]}$ is the scalar equation $\acvs_0\beta_1=\acvs_1$, so
$\beta_1=\acvs_1/\acvs_0=\acf_1$ and
\[
X^{\,1}_2=\beta_1X_1=\frac{\acvs_1}{\acvs_0}X_1=\acf_1X_1 .
\]
\textbf{Based on $X_1,X_2$ (so $n=2$, $h=1$).} We could build $\Gamma_2$ and
solve, but it is quicker to use the model directly. Since
$X_3=\phi_1X_2+\phi_2X_1+\eps_3$ we have
$X_3-(\phi_1X_2+\phi_2X_1)=\eps_3$, and $\eps_3$ (white noise, future) is
uncorrelated with $X_1$ and $X_2$:
\[
\E\big[(X_3-\phi_1X_2-\phi_2X_1)X_1\big]=\E[\eps_3X_1]=0,
\qquad
\E\big[(X_3-\phi_1X_2-\phi_2X_1)X_2\big]=\E[\eps_3X_2]=0 .
\]
Thus the candidate $\phi_1X_2+\phi_2X_1$ has error orthogonal to both
observations, so by the prediction-equations theorem it \emph{is} the best
linear predictor: $X^{\,2}_3=\phi_1X_2+\phi_2X_1$. (This is the general fact: for
an AR$(p)$ with $n\ge p$ observations, the one-step forecast is simply
$\sum_{i=1}^p\phi_iX_{n+1-i}$.)
\end{solution}

\begin{exercise}{Prediction equations with non-zero mean \quad\textnormal{[Sheet 10]}}{p8-ex2}
For a stationary process with $\mu=\E[X_t]$ and predictor
$X^{\,n}_{n+h}=\beta_0+\sum_{j=1}^n\beta_jX_j$ (set $X_0\equiv1$), prove that the
prediction equations $\E[(X_{n+h}-X^{\,n}_{n+h})X_k]=0$, $k=0,\dots,n$, give
$\Gamma_n\beta=\acvs_{[h]}$ and $\beta_0=\mu(1-\beta\Tr\mathbf1_n)$, and that
when $\Gamma_n$ is invertible
$X^{\,n}_{n+h}=\mu+\acvs_{[h]}\Tr\Gamma_n^{-1}(X-\mu\mathbf1_n)$ with MSE
$\acvs_0-\acvs_{[h]}\Tr\Gamma_n^{-1}\acvs_{[h]}$.
\end{exercise}
\begin{solution}
\textbf{(a) MSE in matrix form.} With $X=(X_1,\dots,X_n)\Tr$ and
$\mathbf1_n$ the all-ones vector,
\[
\begin{aligned}
S(\beta_0,\dots,\beta_n)
&=\E\Big[\big(X_{n+h}-\beta_0-\beta\Tr X\big)^2\Big]\\
&=\Var\big(X_{n+h}-\beta\Tr X\big)+\big(\E[X_{n+h}-\beta_0-\beta\Tr X]\big)^2\\
&=\acvs_0-2\beta\Tr\Cov(X_{n+h},X)+\beta\Tr\Var(X)\beta+\big(\mu-\beta_0-\mu\beta\Tr\mathbf1_n\big)^2,
\end{aligned}
\]
a quadratic bounded below by $0$, hence minimised where all
$\partial S/\partial\beta_r=0$.\\
\textbf{(b) Solve for $\beta_0$.} Only the last term depends on $\beta_0$:
\[
\frac{\partial S}{\partial\beta_0}=2\big(\mu-\beta_0-\mu\beta\Tr\mathbf1_n\big)(-1)
=0
\;\Longrightarrow\;
\beta_0=\mu\big(1-\beta\Tr\mathbf1_n\big).
\]
\textbf{(c) Reduce to zero mean.} Substituting $\beta_0$ back,
\[
X^{\,n}_{n+h}=\beta_0+\beta\Tr X=\mu+\beta\Tr(X-\mu\mathbf1_n),
\]
so the predictor depends on the data only through the \emph{centred} vector
$X-\mu\mathbf1_n$. Hence there is no loss of generality in taking $\mu=0$, in
which case $\beta_0=0$ and $X^{\,n}_{n+h}=\beta\Tr X$.\\
\textbf{(d) Normal equations.} For $\mu=0$,
$S(\beta)=\acvs_0-2\beta\Tr\Cov(X_{n+h},X)+\beta\Tr\Var(X)\beta$; setting
$\nabla_\beta S=0$ gives
\[
\Var(X)\,\beta=\Cov(X_{n+h},X),\quad\text{i.e.}\quad \Gamma_n\beta=\acvs_{[h]}.
\]
\textbf{(e) Equivalence to the orthogonality form.} Row $k$ of
$\Gamma_n\beta=\acvs_{[h]}$ reads $\sum_{j=1}^n\beta_j\acvs_{j-k}=\acvs_{n+h-k}$,
which is exactly $\E[(X_{n+h}-\sum_j\beta_jX_j)X_k]=0$ for $k=1,\dots,n$. The
$k=0$ equation (with $X_0=1$) gives $\mu-\beta_0-\sum_j\beta_j\mu=0$, recovering
$\beta_0=\mu(1-\beta\Tr\mathbf1_n)$ from part (b).\\
\textbf{(2a) Closed form.} If $\Gamma_n$ is invertible,
$\beta=\Gamma_n^{-1}\acvs_{[h]}$, so
\[
X^{\,n}_{n+h}=\mu+\acvs_{[h]}\Tr\Gamma_n^{-1}(X-\mu\mathbf1_n),
\]
differing from the lecture's zero-mean formula only by the centring/shift by
$\mu$.\\
\textbf{(2b) MSE.} Since $\E[X_{n+h}-X^{\,n}_{n+h}]=0$,
\[
P^{\,n}_{n+h}=\Var\big(X_{n+h}-\beta\Tr X\big)
=\acvs_0-2\beta\Tr\acvs_{[h]}+\beta\Tr\Gamma_n\beta
=\acvs_0-\acvs_{[h]}\Tr\Gamma_n^{-1}\acvs_{[h]},
\]
identical to the zero-mean case (the mean does not affect forecast accuracy).
\end{solution}

\begin{exercise}{Causal-ARMA predictor via $\psi$-weights \quad\textnormal{[Sheet 10]}}{p8-ex3}
Prove that the best linear predictor $X^{\,n}_{n+h}$ for $X_{n+h}$ in a causal
ARMA process with linear representation
$X_t=\sum_{j=0}^\infty\psi_j\eps_{t-j}$ is
$X^{\,n}_{n+h}=\sum_{j=h}^\infty\psi_j\eps_{n+h-j}$, with MSE
$\sigma^2\sum_{j=0}^{h-1}\psi_j^2$.
\end{exercise}
\begin{solution}
Consider a linear predictor $X^{\,n}_{n+h}=\sum_{i=1}^n\beta_iX_i$. Substituting
the linear representation $X_i=\sum_{j=0}^\infty\psi_j\eps_{i-j}$,
\[
X^{\,n}_{n+h}=\sum_{i=1}^n\beta_i\sum_{j=0}^\infty\psi_j\eps_{i-j}
=\sum_{j=0}^\infty c_j\,\eps_{n-j}
\]
for some coefficients $c_j$; i.e.\ \emph{any} linear predictor based on
$X_1,\dots,X_n$ is a linear combination of innovations up to time $n$ (it
``starts at $\eps_n$''). Now compute the MSE, using
$X_{n+h}=\sum_{j=0}^\infty\psi_j\eps_{n+h-j}$ and splitting the future
innovations $\eps_{n+1},\dots,\eps_{n+h}$ (indices $j=0,\dots,h-1$) from the
rest:
\[
\begin{aligned}
\E\Big[\big(X_{n+h}-X^{\,n}_{n+h}\big)^2\Big]
&=\E\Big[\Big(\sum_{j=0}^\infty\psi_j\eps_{n+h-j}-\sum_{j=0}^\infty c_j\eps_{n-j}\Big)^2\Big]\\
&=\sigma^2\Big(\sum_{j=0}^{h-1}\psi_j^2+\sum_{j=h}^\infty(\psi_j-c_{j-h})^2\Big),
\end{aligned}
\]
where the first sum collects the future innovations (which no past-based
predictor can touch) and the second matches $\psi_j$ against $c_{j-h}$ for the
observable innovations (cross terms vanish because the $\eps$'s are
uncorrelated). The first sum is fixed; the second is minimised --- to zero --- by
choosing $c_{j-h}=\psi_j$, i.e.\ $c_j=\psi_{j+h}$ for $j=0,1,\dots$. The
predictor is then
\[
X^{\,n}_{n+h}=\sum_{j=0}^\infty c_j\eps_{n-j}=\sum_{j=0}^\infty\psi_{j+h}\eps_{n-j}
=\sum_{j=h}^\infty\psi_j\eps_{n+h-j},
\]
and the residual is $X_{n+h}-X^{\,n}_{n+h}=\sum_{j=0}^{h-1}\psi_j\eps_{n+h-j}$,
giving
\[
P^{\,n}_{n+h}=\E\Big[\big(X_{n+h}-X^{\,n}_{n+h}\big)^2\Big]=\sigma^2\sum_{j=0}^{h-1}\psi_j^2 .
\]
\end{solution}

\begin{exercise}{ARIMA(3,1,1) worked truncated forecast \quad\textnormal{[Lecture]}}{p8-ex4}
For the fitted model $(1-\phi_1B-\phi_2B^2-\phi_3B^3)(1-B)y_t=(1+\theta_1B)\eps_t$
(the convention used in the lecture's worked example, with the MA polynomial
written $\Theta(B)=1+\theta_1B$), derive the one-, two- and general $h$-step
truncated forecasts. Then plug in the fitted values $\hat\phi_1=0.0044$,
$\hat\phi_2=0.0916$, $\hat\phi_3=0.3698$, $\hat\theta_1=-0.3921$ and write the
explicit $\hat y_{T+1\mid T}$ and $\hat y_{T+2\mid T}$.
\end{exercise}
\begin{solution}
\textbf{Step 1 --- isolate $y_t$.} Multiply out the AR$\times$difference factor
$(1-\phi_1B-\phi_2B^2-\phi_3B^3)(1-B)$:
\[
1-(1+\phi_1)B+(\phi_1-\phi_2)B^2+(\phi_2-\phi_3)B^3+\phi_3B^4,
\]
so the model is
\[
y_t-(1+\phi_1)y_{t-1}+(\phi_1-\phi_2)y_{t-2}+(\phi_2-\phi_3)y_{t-3}+\phi_3y_{t-4}
=\eps_t+\theta_1\eps_{t-1}.
\]
Moving everything except $y_t$ to the right,
\[
y_t=\underbrace{(1+\phi_1)y_{t-1}-(\phi_1-\phi_2)y_{t-2}-(\phi_2-\phi_3)y_{t-3}-\phi_3y_{t-4}}_{\text{AR part}}
+\underbrace{\eps_t+\theta_1\eps_{t-1}}_{\text{MA part}} .
\]
(We use the lecture's convention $\Theta(B)=1+\theta_1B$, so the MA term is
$\eps_t+\theta_1\eps_{t-1}$ and the forecast picks up $+\theta_1\hat\eps$.
Beware: in the earlier general theory the MA polynomial was written
$\Theta(B)=1-\theta_1B$; that sign convention would replace every $+\theta_1$
below by $-\theta_1$. We follow the worked example exactly so the numbers match.)
\\[3pt]
\textbf{Step 2--3 --- shift and substitute.} For $h=1$ (replace $t$ by $n+1$):
the only error terms are $\eps_{n+1}\to0$ and $\eps_n\to\hat\eps_n$ (last
residual). Hence
\[
y^{\,n}_{n+1}=(1+\phi_1)y_n-(\phi_1-\phi_2)y_{n-1}-(\phi_2-\phi_3)y_{n-2}-\phi_3y_{n-3}+\theta_1\hat\eps_n .
\]
For $h>1$ all error terms are future ($\to0$) and the future observations are
replaced by forecasts:
\[
y^{\,n}_{n+h}=(1+\phi_1)y^{\,n}_{n+h-1}-(\phi_1-\phi_2)y^{\,n}_{n+h-2}-(\phi_2-\phi_3)y^{\,n}_{n+h-3}-\phi_3y^{\,n}_{n+h-4},
\]
with the convention $y^{\,n}_s=y_s$ for $s\le n$.\\[3pt]
\textbf{Numbers.} Using
$1+\hat\phi_1=1.0044$, $\hat\phi_1-\hat\phi_2=0.0044-0.0916=-0.0872$,
$\hat\phi_2-\hat\phi_3=0.0916-0.3698=-0.2782$, $\hat\phi_3=0.3698$,
$\hat\theta_1=-0.3921$, the one-step forecast (writing $e_T$ for the last
residual) is
\[
\hat y_{T+1\mid T}=1.0044\,y_T+0.0872\,y_{T-1}+0.2782\,y_{T-2}-0.3698\,y_{T-3}+\hat\theta_1\,e_T,
\]
i.e.\ with $\hat\theta_1=-0.3921$,
\[
\hat y_{T+1\mid T}=1.0044\,y_T+0.0872\,y_{T-1}+0.2782\,y_{T-2}-0.3698\,y_{T-3}-0.3921\,e_T .
\]
(The two middle signs are positive because $-(\hat\phi_1-\hat\phi_2)=+0.0872$ and
$-(\hat\phi_2-\hat\phi_3)=+0.2782$.) The two-step forecast replaces $y_{T+1}$ by
$\hat y_{T+1\mid T}$ and sets all errors to $0$:
\[
\hat y_{T+2\mid T}=1.0044\,\hat y_{T+1\mid T}+0.0872\,y_T+0.2782\,y_{T-1}-0.3698\,y_{T-2}.
\]
Continue likewise for $h=3,4,\dots$; each step uses previously computed
forecasts. Because $d=1$, the prediction MSE grows without bound as $h$
increases.
\end{solution}

\begin{exercise}{One- and two-step forecast of an MA(2) with interval \quad\textnormal{[New]}}{p8-ex5}
Let $X_t=\eps_t-\theta_1\eps_{t-1}-\theta_2\eps_{t-2}$ with
$\eps_t\iid\mathcal N(0,\sigma^2)$ (invertible, zero mean). Give
$X^{\,n}_{n+1}$, $X^{\,n}_{n+2}$, $X^{\,n}_{n+3}$, their forecast-error
variances, and a $95\%$ interval for $X_{n+2}$.
\end{exercise}
\begin{solution}
Here $\psi_0=1,\ \psi_1=-\theta_1,\ \psi_2=-\theta_2$ and $\psi_j=0$ for
$j\ge3$. By the $\psi$-weight theorem, set future innovations to $0$ in
$X_{n+h}=\eps_{n+h}-\theta_1\eps_{n+h-1}-\theta_2\eps_{n+h-2}$:
\[
\begin{aligned}
X^{\,n}_{n+1}&=-\theta_1\hat\eps_n-\theta_2\hat\eps_{n-1},\\
X^{\,n}_{n+2}&=-\theta_2\hat\eps_n,\\
X^{\,n}_{n+3}&=0\quad(=\mu,\ \text{the mean; an MA(2) is forecast as the mean for }h>2).
\end{aligned}
\]
Forecast-error variances $\Var(e_n(h))=\sigma^2\sum_{j=0}^{h-1}\psi_j^2$:
\[
\Var(e_n(1))=\sigma^2,\quad
\Var(e_n(2))=\sigma^2(1+\theta_1^2),\quad
\Var(e_n(3))=\sigma^2(1+\theta_1^2+\theta_2^2)=\acvs_0 .
\]
For $h\ge3$ the variance saturates at $\acvs_0$, as it must. A $95\%$ interval
for $X_{n+2}$ is
\[
-\theta_2\hat\eps_n\;\pm\;1.96\,\sigma\sqrt{1+\theta_1^2}.
\]
\end{solution}

\begin{exercise}{Two-step AR(1) forecast variance and interval \quad\textnormal{[New]}}{p8-ex6}
For the Gaussian AR(1) $X_t=\phi X_{t-1}+\eps_t$ ($\abs\phi<1$,
$\eps_t\iid\mathcal N(0,\sigma^2)$), compute $X^{\,n}_{n+1},X^{\,n}_{n+2}$,
their error variances, and verify $\Var(e_n(h))\to\acvs_0$.
\end{exercise}
\begin{solution}
The $MA(\infty)$ weights of a causal AR(1) are $\psi_j=\phi^j$. Point forecasts
(Example~\ref{exmp:p8-ar1}): $X^{\,n}_{n+1}=\phi X_n$,
$X^{\,n}_{n+2}=\phi^2X_n$, and generally $X^{\,n}_{n+h}=\phi^hX_n$.
Forecast-error variances:
\[
\Var(e_n(1))=\sigma^2\psi_0^2=\sigma^2,
\qquad
\Var(e_n(2))=\sigma^2(\psi_0^2+\psi_1^2)=\sigma^2(1+\phi^2),
\]
and in general
$\Var(e_n(h))=\sigma^2\sum_{j=0}^{h-1}\phi^{2j}=\sigma^2\dfrac{1-\phi^{2h}}{1-\phi^2}$.
As $h\to\infty$ this tends to $\dfrac{\sigma^2}{1-\phi^2}=\acvs_0$, the marginal
variance of an AR(1) --- confirming the saturation. A $95\%$ interval for
$X_{n+2}$ is $\phi^2X_n\pm1.96\,\sigma\sqrt{1+\phi^2}$.
\end{solution}

\begin{exercise}{Re-integrating a random walk with drift \quad\textnormal{[New]}}{p8-ex7}
Let $X_t$ be ARIMA$(0,1,0)$ with drift: $Z_t=X_t-X_{t-1}=\alpha+\eps_t$,
$\eps_t\iid\mathcal N(0,\sigma^2)$. Find $X^{\,n}_{n+h}$ and $\Var(e_n(h))$, and
show the MSE diverges.
\end{exercise}
\begin{solution}
The differenced series $Z_t$ is white noise plus a constant mean
$\E[Z_t]=\alpha$, so $Z^{\,n}_{n+h}=\alpha$ for all $h$. Re-integrating
(Proposition on integrated processes),
\[
X^{\,n}_{n+h}=X_n+\sum_{j=0}^{h-1}Z^{\,n}_{n+h-j}=X_n+\alpha h,
\]
a straight drift line from the last observation. The forecast error is
$e_n(h)=X_{n+h}-X^{\,n}_{n+h}=\sum_{j=1}^h\eps_{n+j}$ (the accumulated future
innovations), so
\[
\Var(e_n(h))=h\,\sigma^2\;\xrightarrow[h\to\infty]{}\;\infty .
\]
Thus $P^{\,n}_{n+h}=h\sigma^2$ grows linearly and the prediction interval
$X_n+\alpha h\pm1.96\,\sigma\sqrt h$ fans out without bound --- the hallmark of an
integrated process (no mean reversion).
\end{solution}
